\section{Result 3: part-way through training, the scores look like your prediction}
\label{sec:equal}

While we were planning you suggested another way to compare: give every model the same amount
of text and then see how well it does, without waiting for each one to finish. We left that
for after the main experiment. It needs no new training. Every run was graded as it went
along, so I can read off what each model scored after the same number of sets, which is the
same amount of text in every language.

\begin{table}[H]
\centering\small
\begin{tabular}{@{}rrrrr@{\qquad}rrrr@{}}
\toprule
 & \multicolumn{4}{c}{After 32 sets} & \multicolumn{4}{c}{After 64 sets} \\
Width & 16 & 64 & 256 & 1{,}024 & 16 & 64 & 256 & 1{,}024 \\
\midrule
8 & 69 & 37 & 5 & 0 & 100 & 74 & 17 & 2 \\
16 & 100 & 97 & 59 & 13 & 100 & 100 & 100 & 51 \\
32 & 100 & 100 & 81 & 32 & 100 & 100 & 100 & 75 \\
\bottomrule
\end{tabular}

\caption{Facts known (0 to 100) after 32 sets and after 64 sets of text, by width and by the
number of animals in the language. First round, average of the three repeats. A run that was
already solved by then is entered with its final score. I show 32 and 64 sets because earlier
most entries are near 0 and later most are 100.}
\label{tab:equal}
\end{table}

\textbf{Part-way through training, the scores look like your prediction} (table
\ref{tab:equal}). Take width 16 after 32 sets. It scores 100 on 16 animals, 59 on 256 and 13 on
1,024: the more words, the lower the score. And a wider model scores higher. On 256 animals
after 32 sets, width 8 scores 5, width 16 scores 59 and width 32 scores 81. The unequal-use
round shows the same pattern: after 32 sets on 1,024 animals the three widths score 19, 79 and
99.

\textbf{But the cause is not the one your idea names.} Your idea is about room: a model that
is full cannot take in more words however long it is trained. Nothing here shows a full
model. With more text every one of these models went on to a score of 100 (section
\ref{sec:round1}). What the table shows is speed. A larger language takes more text to learn,
and a wider model learns it from less.

\textbf{A larger language needs more text in all, but less text for each fact.} In all: width
8 needed 512 sets for 1,024 animals against 64 sets for 16 (the middle one of its three repeats
each time), 8 times the text for 64 times the facts. For each fact: one set shows each fact of
the smallest language 960 times and each fact of the largest only 15 times. So by the time
width 8 was solved, each fact had come up about 61,000 times in the small language and about
7,700 times in the large one.

\textbf{How much weight to give this.} What the comparison shows is true and useful: with a
limited amount of text, the same model knows less of a larger language, and a wider model knows
more of it. What it cannot show is why. A model that has run out of room and a model that only
needs more text look the same at one moment in training. Here we know it was the second,
because with more text every model reached 100. And although you proposed the comparison, we
fixed no rule for reading it, and I chose the two moments in the table after seeing the
results. So it is an observation, not a test.
